
所需的仓西结构现已构造完成。
下面为实现算子 $\mathfrak p$ 的构造，
在模空间
$\mathcal M^{\mathrm{main}}_{k,1}(L;\beta)$
（$k=0,1,2,\ldots$）上
构造一套多值截面（扰动）。
这也就是给出
定理~\ref{thm:3-8-9}
和引理~\ref{lem:3-8-25}
证明的细节。

须注意，在扰动
$\mathcal M^{\mathrm{main}}_{k,1}(L;\beta)$
之前，
我们已完成
$\mathcal M^{\mathrm{main}}_{k,0}
(L;\beta;P_1,\ldots,P_k)$ 上
扰动的构造，
也选好了奇异单形可数集 $\mathcal X(L)$。
还须固定
$\mathcal M_k(M;\widetilde\beta)$
上的扰动，使之与球面气泡相容。
记
\begin{equation}
\begin{aligned}
&\mathcal M^{\mathrm{main}}_{k,1}
(L;\beta;P_1,\ldots,P_k)\\
&\qquad=
\mathcal M^{\mathrm{main}}_{k,1}(L;\beta)
\times_{L^k}(P_1\times\cdots\times P_k).
\end{aligned}
\tag{7.4.7}\label{eq:7-4-7}
\end{equation}
这个空间带有仓西结构。
当 $k=0$ 时，
须将 $\mathcal M^{\mathrm{main}}_{0,1}(L;\beta)$
紧化为
\begin{equation}
\mathcal M^{\mathrm{main}}_{0,1}(L;\beta)
\cup
\bigcup_{\substack{\widetilde\beta:\\
i_*(\widetilde\beta)=\beta}}
\bigl(\mathcal M_2(M;\widetilde\beta)
\times_M L\bigr).
\tag{7.4.8}\label{eq:7-4-8}
\end{equation}
以下仍以
$\mathcal M^{\mathrm{main}}_{0,1}(L;\beta)$
表示这一紧化，以简化记号。

我们对 $\omega[\beta]$ 归纳，
在
$\mathcal M^{\mathrm{main}}_{k,1}
(L;\beta;P_1,\ldots,P_k)$
以及
$\mathcal M^{\mathrm{main}}_{0,1}(L;\beta)$
上构造多值截面系统。
前者的多值截面可能依赖于
$k'>k$、但 $\omega[\beta']<\omega[\beta]$
时
$\mathcal M^{\mathrm{main}}_{k',1}
(L;\beta';P'_1,\ldots,P'_{k'})$
上多值截面的选择；
因此不同 $k$ 的构造必须同时进行。

\begin{remark}{7.4.9}\label{rem:7-4-9}
为
$\mathcal M^{\mathrm{main}}_{k,1}
(L;\beta;P_1,\ldots,P_k)$
构造多值截面时，
并不需要第~\ref{sec:7-2} 节
那样精细的归纳。
关键区别是：
研究 $\mathfrak m$ 时，
还需把 $\mathfrak m$ 作用于
$\mathfrak m_k(P_1,\ldots,P_k)$；
而 $\mathfrak p_k(P_1,\ldots,P_k)$
是环境辛流形中的奇异单形，
不再需要将 $\mathfrak m$ 或
$\mathfrak p$ 作用于
奇异链 $\mathfrak p_k(P_1,\ldots,P_k)$。
\end{remark}

规定 $(\beta,k)>(\beta',k')$
当且仅当下列条件之一成立：
\par\noindent\textup{(7.4.10.1)}\label{cond:7-4-10-1}\quad
$\omega[\beta]>\omega[\beta']$；
\par\noindent\textup{(7.4.10.2)}\label{cond:7-4-10-2}\quad
$\omega[\beta]=\omega[\beta']$ 且 $k>k'$。
依此顺序归纳构造
$\mathcal M^{\mathrm{main}}_{k,1}
(L;\beta;P_1,\ldots,P_k)$
上的多值截面。

注意
$\mathcal M^{\mathrm{main}}_{0,1}(L;\beta_0)
=\varnothing$：
其中没有稳定映射，
且
$\mathcal M_2(M;\widetilde\beta_0)
=\varnothing$。
故 $\mathfrak p_{0,\beta_0}=0$。
也不必扰动
$\mathcal M^{\mathrm{main}}_{1,1}(L;\beta_0)$；
具体地，对任意 $P$ 置
$\mathcal M_{1,1}(L;\beta_0;P)=P$。

现假设对一切
$(\beta',k')<(\beta,k)$
的情形，
或在 $(\beta',k')=(\beta,k)$
而
$\sum_i\dim P'_i<\sum_i\dim P_i$
的情形，
已构造
$\mathcal M^{\mathrm{main}}_{k',1}
(L;\beta';P'_1,\ldots,P'_{k'})$
上的扰动（多值截面）。
在这些归纳假设下，
构造
$\mathcal M^{\mathrm{main}}_{k,1}
(L;\beta;P_1,\ldots,P_k)$
上的扰动。

考察其边界
$\partial\mathcal M^{\mathrm{main}}_{k,1}
(L;\beta;P_1,\ldots,P_k)$，
可分为下列几类；
见图~\ref{fig:7-4-2}、
\ref{fig:7-4-3}。

\par\noindent\textup{(7.4.11.1)}\label{cond:7-4-11-1}\quad
由输入单形的边界产生：
\[
\mathcal M^{\mathrm{main}}_{k,1}
(L;\beta;P_1,\ldots,\partial P_i,\ldots,P_k).
\]

\par\noindent\textup{(7.4.11.2)}\label{cond:7-4-11-2}\quad
由图~\ref{fig:7-4-2} 的圆盘分裂产生：
\[
\mathcal M^{\mathrm{main}}_{k-j+i,1}
\bigl(L;\beta-\beta';
P_1,\ldots,P_{i-1},
Q_{i,j}(\beta'),P_{j+1},\ldots,P_k\bigr),
\]
其中 $1\leq i\leq j+1\leq k+1$，且
\[
Q_{i,j}(\beta')
=\operatorname{ev}_{0*}
\bigl(\mathcal M^{\mathrm{main}}_{j-i+2,0}
(L;\beta';P_i,\ldots,P_j)\bigr).
\]

\par\noindent\textup{(7.4.11.3)}\label{cond:7-4-11-3}\quad
由图~\ref{fig:7-4-3} 的另一种圆盘分裂产生：
\[
\mathcal M^{\mathrm{main}}_{j-i+2,1}
\bigl(L;\beta-\beta';
P_i,\ldots,P_j,R_{j+1,i-1}(\beta')\bigr),
\]
其中 $2\leq i\leq j+1\leq k+1$，且
\[
R_{j+1,i-1}(\beta')
=\operatorname{ev}_{0*}
\bigl(\mathcal M^{\mathrm{main}}_{k-j+i,0}
(L;\beta';P_{j+1},\ldots,P_k,
P_1,\ldots,P_{i-1})\bigr).
\]

\par\noindent\textup{(7.4.11.4)}\label{cond:7-4-11-4}\quad
当 $k=0$ 时，还出现
$\mathcal M_2(M;\widetilde\beta)\times_M L$。

在扰动之前，
这里有循环对称性\liftIndexRangeStart{I385}：
$\mathbb Z_k$ 作用由
\[
\operatorname{cyc}\colon
\mathcal M^{\mathrm{main}}_{k,1}
(L;\beta;P_1,\ldots,P_k)
\longrightarrow
\mathcal M^{\mathrm{main}}_{k,1}
(L;\beta;P_2,\ldots,P_k,P_1)
\]
生成。
我们的目标是构造保持此循环对称性的
扰动（多值截面）。
式~(7.4.11) 中各模空间的扰动
已由归纳假设给出。
图~\ref{fig:7-4-3} 的右侧圆盘
按边界型别标为 $\beta-\beta'$
\translatornote{源书图 7.4.3
将右侧圆盘标为 $\beta$；
但所对应的边界型别 (7.4.11.3)
明确要求该分支为 $\beta-\beta'$，
故重绘图采用后者。
见疑误 A-0558。}。
\par\medskip\noindent\begin{minipage}{\textwidth}
\centering
\setlength{\unitlength}{0.65mm}
\begin{picture}(137,91)
\linethickness{0.65pt}
{\color{gray!55}
\put(38,53){\circle{47}}
\put(88,53){\circle{65}}}
\put(38,53){\circle{42}}
\put(88,53){\circle{58}}
\put(59,53){\circle*{1.5}}
\put(103,53){\circle*{1.5}}
\put(100,48){$\scriptstyle z_0$}
\put(30,52){$\beta'$}
\put(79,52){$\beta-\beta'$}
\put(20,69){\circle*{1.2}}
\put(12,70){$\scriptstyle P_i$}
\put(49,37){\circle*{1.2}}
\put(45,26){$\scriptstyle P_j$}
\put(72,78){\circle*{1.2}}
\put(59,79){$\scriptstyle P_{i-1}$}
\put(73,29){\circle*{1.2}}
\put(60,20){$\scriptstyle P_{j+1}$}
\put(108,72){\circle*{1.2}}
\put(111,73){$\scriptstyle P_1$}
\put(110,33){\circle*{1.2}}
\put(113,28){$\scriptstyle P_k$}
\put(59,5){\vector(-1,0){23}}
\put(59,5){\line(0,1){10}}
\put(20,1){$\scriptstyle Q_{i,j}(\beta')$}
\end{picture}
\par\smallskip
{\renewcommand{\thefigure}{7.4.2}\refstepcounter{figure}
\normalfont\bfseries 图~\thefigure\label{fig:7-4-2}\quad
\normalfont 情形 (7.4.11.2)}
\end{minipage}\par\medskip

\par\medskip\noindent\begin{minipage}{\textwidth}
\centering
\setlength{\unitlength}{0.65mm}
\begin{picture}(137,90)
\linethickness{0.65pt}
{\color{gray!55}
\put(42,52){\circle{58}}
\put(94,52){\circle{58}}}
\put(42,52){\circle{52}}
\put(94,52){\circle{52}}
\put(68,52){\circle*{1.5}}
\put(25,51){\circle*{1.5}}
\put(19,44){$\scriptstyle z_0$}
\put(37,50){$\beta'$}
\put(83,50){$\beta-\beta'$}
\put(17,65){\circle*{1.2}}
\put(5,66){$\scriptstyle P_k$}
\put(16,50){\circle*{1.2}}
\put(1,47){$\scriptstyle P_1$}
\put(61,69){\circle*{1.2}}
\put(53,71){$\scriptstyle P_{j+1}$}
\put(60,35){\circle*{1.2}}
\put(49,25){$\scriptstyle P_{i-1}$}
\put(78,72){\circle*{1.2}}
\put(77,76){$\scriptstyle P_j$}
\put(94,78){\circle*{1.2}}
\put(95,79){$\scriptstyle P$}
\put(79,30){\circle*{1.2}}
\put(79,19){$\scriptstyle P_i$}
\put(68,4){\vector(-1,0){23}}
\put(68,4){\line(0,1){11}}
\put(10,0){$\scriptstyle R_{j+1,i-1}(\beta')$}
\end{picture}
\par\smallskip
{\renewcommand{\thefigure}{7.4.3}\refstepcounter{figure}
\normalfont\bfseries 图~\thefigure\label{fig:7-4-3}\quad
\normalfont 情形 (7.4.11.3)}
\end{minipage}\par\medskip

\liftIndexRangeEnd{I385}
